\section{Conclusion}
\label{sec:conclusion}

This paper has provided the normative framework that situates the six L-series
metrics in their broader democratic theory context. The proportionality-majoritarianism
debate is not resolvable by empirical analysis---it is a normative question about
what representation should look like, and courts applying state constitutional
standards must resolve it as a matter of law, not measurement. Expert witnesses'
role is to provide the empirical foundation under whichever normative standard
the court applies.

The central empirical findings are robust. \textsc{Bisect} maps achieve
$\text{LSq} \in [0.03, 0.08]$ across four states, producing a majoritarian bonus
of $1.3$--$1.5$---within the historical Anglo-American average for single-member
plurality systems. Enacted maps achieve $\text{LSq} \in [0.10, 0.18]$, producing
a majoritarian bonus of $1.9$--$2.2$---substantially above the structural baseline,
reflecting redistricting choices that lock in the majority party's advantage
beyond what competitive district outcomes would produce.

The normative contribution is the reframing of the proportionality defense. When
defendants argue that ``proportionality is not required,'' the correct expert
witness response is: (1) agree that strict proportionality ($\text{LSq} = 0$)
is not required at the federal level or under most state constitutions; (2) show
that the neutral-criteria baseline ($\text{LSq} \approx 0.03$--$0.08$) is the
minimum disproportionality that any compact, population-equal redistricting plan
must produce for this state; (3) show that the enacted plan's disproportionality
($\text{LSq} \approx 0.10$--$0.18$) is substantially above this structural floor;
and (4) invite the court to determine whether the gap above the neutral baseline
is permissible under the applicable state constitutional standard. This reframing
converts a defense that concedes everything the defendant needs (``proportionality
not required'') into a question that courts applying ``free and equal elections''
standards can readily evaluate: is the mapmaker's contribution to disproportionality
consistent with a constitution that requires elections to be free from deliberate
partisan manipulation?

\textsc{Bisect} algorithmic maps satisfy both the proportionality standard (to
the extent single-member districts can) and the majoritarian standard (they
produce a competitive majority bonus without structurally locking in partisan
advantage) more closely than any enacted plan in the four states studied.
This convergence is the paper's ultimate contribution to partisan fairness expert
testimony: regardless of which normative framework applies, the algorithmic maps
provide a superior reference distribution.

\bigskip

\noindent\textit{The author thanks the redistricting research team for data
access and algorithmic development. All errors are the author's own.}
